\subsection{Fiber products}

We will show that the category of schemes has fiber products (pullbacks).
We do this in steps.

\blemm[name={pullbackaffine}]

    Suppose $\alpha\colon X\to Z$ and $\beta\colon Y\to Z$ are morphisms of affine schemes.
    Then their fiber product $X\times_ZY$ exists (in the category of schemes) and is affine.

\elemm

\bproof

    Let $X=\spec A$, $Y=\spec B$, $Z=\spec C$, then their fiber product is a representation of
    $$ \Sch(\bul,\spec A)\times_{\Sch(\bul,\spec C)}\Sch(\bul,\spec B) , $$
    and following the spec-global sections functor we see that
    $$ \eqalign{
        \Sch(\bul,\spec A)\times_{\Sch(\bul,\spec C)}\Sch(\bul,\spec B) &\cong
            \Ring(A,\Gamma\bul)\times_{\Ring(C,\Gamma\bul)}\Ring(B,\Gamma\bul)\cr
        &\cong\Ring(A\otimes_CB,\Gamma\bul)\cr
        &\cong \Sch(\bul,\spec(A\otimes_CB)) .
    } $$
    So that the fiber product is $\spec(A\otimes_CB)$ which is affine.
    \qed

\eproof

\bnote

    It follows that $\bA^1_k\times_k\bA^1_k\cong\bA^2_k$, because $k[x]\otimes_kk[y]\cong k[x,y]$
    via the obvious map $p(x)\otimes q(y)\mapsto p(x)q(y)$.
    But as sets it is not true that $\abs{\bA^2_k}\cong\abs{\bA^1_k}\times\abs{\bA^1_k}$, as there are
    prime ideals of $\bA^2_k$ not of the form $(x-a,y-b)$ (i.e.\ which are not maximal) or $(x-a)$ or
    $(y-b)$.
    So the underlying set of a fiber product of schemes is not the fiber product of underlying sets.

\enote

\blemm[name={pullbackopen}]

    Suppose $\iota\colon U\inj Z$ is an open embedding and $\rho\colon Y\to Z$ is a morphism.
    Then $U\times_ZY$ is $(\rho^{-1}U,\O_Y|_{\rho^{-1}U})$.
    In particular the fiber product of an open embedding is an open embedding:

    \diagram{
        U\times_ZY&Y\cr
        U&Z
    }{
        \diagarrow{from={1,1}, to={1,2}, left cap={(}, text=\text{open}, y distance=.15cm}
        \diagarrow{from={1,1}, to={2,1}}
        \diagarrow{from={2,1}, to={2,2}, left cap={(}, text=\text{open}, y distance=.15cm}
        \diagarrow{from={1,2}, to={2,2}}
    }

\elemm

\bproof

    Suppose we have the following commutative square
    \commsq{W&Y\cr U&Z}
        {\alpha}{\iota}{\beta}{\rho}
    Then as topological maps we must have that $\alpha(W)\subseteq\rho^{-1}U$.
    And so restricting $\alpha$'s codomain to the open $\rho^{-1}U$ gives us the desired map
    $\bar\alpha\colon W\to\rho^{-1}U$ making the diagram commute.
    \qed

\eproof

\blemm[name={pullbackalmostaffine}]

    Suppose $\alpha\colon X\to Z$ is a morphism of affine schemes, and $\beta\colon Y\to Z$
    factors as $Y\xinj\iota Y'\to Z$ where $\iota$ is an open embedding and $Y'$ affine.
    Then $X\times_ZY$ exists.

\elemm

\bproof

    Consider the diagram below:
    \diagram{
        W&W'&X\cr
        Y&Y'&Z
    }{
        \diagarrow{from={1,1}, to={1,2}, left cap={(}, text=\text{open}, y distance=.15cm}
        \diagarrow{from={1,1}, to={2,1}}
        \diagarrow{from={2,1}, to={2,2}, left cap={(}, text=\text{open}, y distance=.15cm}
        \diagarrow{from={1,2}, to={2,2}}
        \diagarrow{from={1,2}, to={1,3}}
        \diagarrow{from={2,2}, to={2,3}}
        \diagarrow{from={1,3}, to={2,3}}
    }
    The right square has a pullback (because the schemes are affine), and the left square has a pullback
    (because open embeddings have pullbacks).
    Thus the rectangle as a whole is a pullback.
    \qed

\eproof

\blemm[name={pullbackyarb}]

    Suppose $X,Z$ are affine and $Y$ is arbitrary.
    Then $X\times_ZY$ exists.

\elemm

\bproof

    First assume that $Y$ has a cover by two affine open subsets $Y_1,Y_2$.
    Let $Y_{12}=Y_1\cap Y_2$, so that by \refmath[lemma]{pullbackaffine} $W_i=X\times_ZY_i$ exists
    and so does $W_{12}=X\times_ZY_{12}$ by \refmath[lemma]{pullbackalmostaffine}.
    Furthermore by the proof of \refmath[lemma]{pullbackalmostaffine} the canonical map
    $X\times_ZY_{12}\inj X\times_ZY_i$ is an open embedding (it is the map $W\inj W'$ in the proof).
    And so we can glue $W_1$ and $W_2$ along $W_{12}$, call this scheme $W$.

    There are obvious maps $W\xtos{\alpha'}Y$ and $W\xtos{\beta'}X$.
    We want to show that $W=X\times_ZY$.
    So suppose that we have $\alpha''\colon V\to X$ and $\beta''\colon V\to Y$ which commute with
    $\alpha,\beta$.
    That is, we must fill in the below diagram with a $\gamma$:

    \diagram{
        V\cr
        &W&Y\cr
        &X&Z
    }{
        \diagarrow{from={1,1}, to={2,3}, text=\beta'', x distance=.15cm}
        \diagarrow{from={1,1}, to={3,2}, text=\alpha'', x distance=-.15cm}
        \diagarrow{from={1,1}, to={2,2}, text=\gamma, x distance=-.15cm, dashed}
        \diagarrow{from={2,2}, to={2,3}, text=\alpha', y distance=.15cm}
        \diagarrow{from={2,2}, to={3,2}, text=\beta', x distance=-.15cm}
        \diagarrow{from={2,3}, to={3,3}, text=\beta, x distance=.15cm}
        \diagarrow{from={3,2}, to={3,3}, text=\alpha, y distance=.15cm}
    }

    Define $V_i=(\beta'')^{-1}Y_i$ and $V_{12}=(\beta'')^{-1}Y_{12}=V_1\cap V_2$.
    Then for each $i$ there is a unique map $V_i\to W_i$ fitting into the pullback ($W_i=Y_i\times_ZX$).
    Thus a unique $\gamma_i\colon V_i\to W$ fitting into the diagram above (with $V_i$ for $V$).
    And similarly there is a unique map $\gamma_{12}\colon V_{12}\to W$ fitting into the diagram.
    By uniqueness, $\gamma_1$ and $\gamma_2$ must agree on $V_{12}$ with $\gamma_{12}$, and thus there
    is a unique gluing $\gamma\colon V\to W$.

    In the general case, cover $Y$ with an open cover of affine open subsets $Y_i$.
    As before $Y_{ij}=Y_i\cap Y_j$ has a pullback, so let
    $$ W_i = Y_i\times_ZX,\qquad W_{ij} = Y_{ij}\times_ZX . $$
    Then we may glue the $W_i$ along $W_{ij}$ together to form a scheme $W$.

    Now given $\alpha''\colon V\to X$ and $\beta''\colon V\to Y$ as before, define
    $V_i=(\beta'')^{-1}Y_i$ and $V_{ij}=(\beta'')^{-1}Y_{ij}=V_i\cap V_j$.
    As before for each $i$ there is a unique $V_i\to W_i$ fitting into the diagram, giving us a unique
    $\gamma_i\colon V_i\to W$.
    And there are unique $\gamma_{ij}\colon V_{ij}\to W$, so $\gamma_i$ and $\gamma_j$ must agree on
    $V_i\cap V_j$ with $\gamma_{ij}$.
    Thus they uniquely glue together to give $\gamma$.
    \qed

\eproof

\blemm[name={pullbackxyarb}]

    If $Z$ is affine and $X,Y$ are arbitrary then $X\times_ZY$ exists.

\elemm

\bproof

    Consider the above lemma and its proof.
    In it we know that $W_i=X\times_ZY_i$ exists because $Y_i$ is affine (this is what the lemma proves).
    And $W_{ij}=X\times Y_{ij}$ exists by \refmath[lemma]{pullbackopen}.
    Thus the proof of the lemma works again (now that we know the lemma is true) for general $X$ as well.
    \qed

\eproof

\blemm[name={pullbackxyarbzalmostaffine}]

    Suppose $X,Y$ are abitrary and $Z$ is an open subscheme of an affine scheme $Z'$.
    Then $X\times_ZY$ exists.

\elemm

\bproof

    In general if $Z\to Z'$ is monic then $X\times_ZY$ and $X\times_{Z'}Y$ satisfy the same universal property,
    so one exists iff the other does.
    So this lemma follows from the previous one.
    \qed

\eproof

\bthrm[name={fiberproductsexist}]

    For arbitrary schemes $X,Y,Z$, the fiber product $X\times_ZY$ exists.

\ethrm

\bproof

    Cover $Z$ with affine open subschemes $Z_i$, and let $X_i=\alpha^{-1}Z_i$ and $Y_i=\beta^{-1}Z_i$.
    Further define $Z_{ij}=Z_i\cap Z_j$ and $X_{ij}=\alpha^{-1}Z_{ij}=X_i\cap X_j$ and similar for $Y$.
    Then by \refmath[lemma]{pullbackxyarb} $W_iX_i\times_{Z_i}Y_i$ exists for all $i$, and by
    \refmath[lemma]{pullbackxyarbzalmostaffine} $W_{ij}=X_{ij}\times_{Z_{ij}}Y_{ij}$ exists for all $ij$.

    Note that as before $X_i\times_{Z_i}Y_i=X_i\times_ZY_i$, and furthermore we can see that
    $X_i\times_ZY_i=X\times_ZY_i$.
    This is because they both satisfy the same universal property: any $V$ with
    $\alpha''\colon V\to X$ and $\beta''\colon V\to Y_i$ must have that $\alpha\alpha''=\beta\beta''$
    which means that $\alpha\alpha''$ lies in $Z_i$ and thus $\alpha''$ lies in $X_i$.
    So we may identify $W_i=X\times_ZY_i$.
    And similarly $W_{ij}=X\times_ZY_{ij}$, and then we may glue the $W_i$s along the $W_{ij}$s, and for
    the same reason as in \refmath[lemma]{pullbackyarb} this gives us the fiber product
    $X\times_ZY$.
    \qed

\eproof

\subsubsection{Representable functors}

In this section we will give a more sophisticated perspective on fiber products and provide an alternative
proof as to their existence.
We assume only \refmath[lemma]{pullbackaffine} and \refmath[lemma]{pullbackopen}.

Recall that a (contravariant) functor is {\emphcolor representable} if it is naturally isomorphic to
$\hom(\bul,X)$.
And for a diagram $D\colon I\to\C$, its limit exists iff the limit of $\y\circ D\colon I\to\Psh(\C)$ is
representable (here $\y$ is the Yoneda embedding, $\y(X)=\hom(\bul,X)$).
Note that $\y\circ D$'s limit exists because $\Psh(\C)$ is complete.

Writing $h_X=\hom(\bul,X)$, we see that $X\times_ZY$ exists iff $h_X\times_{h_Z}h_Y$ is representable.
So we will give an alternative perspective on the above proof in light of this.

Note that if $U_i$ is an open cover of $Y$, a morphism $Y\to X$ is determined by its restrictions
$U_i\to X$.
And given morphisms $U_i\to X$ which agree on $U_i\cap U_j$, they in turn glue to give $Y\to X$.
Thus we have an equalizer diagram
$$ h_X(Y) \to \prod h_X(U_i) \tto \prod h_X(U_i\cap U_j) , $$
which is to say that $h_X$ forms a sheaf on every scheme $Y$.
In general if $F$ is a contravariant functor $\Sch^\op\to\Set$ which forms a sheaf on every scheme $Y$
(i.e.\ $F|_{\open(Y)}$ is a sheaf), then $F$ is called a {\emphcolor Zariski sheaf}.
So if $F$ is representable it is a Zariski sheaf.

\blemm[name={fiberzariski}]

    For schemes $X,Y,Z$, $h_X\times_{h_Z}h_Y$ is a Zariski sheaf.

\elemm

\bproof

    Let $\alpha\colon X\to Z$ and $\beta\colon Y\to Z$ so that for any scheme $W$
    $$ h_X(W)\times_{h_Z(W)}h_Y(W) = \set{(f\colon W\to X,g\colon W\to Y)}[\alpha f=\beta g] . $$
    Now suppose $U_i$ cover $W$ and we are given $f_i\colon U_i\to X$ and $g_i\colon U_i\to Y$ such that
    $\alpha f_i=\beta g_i$ and $f_i|_{U_i\cap U_j}=f_j|_{U_i\cap U_j}$ and similar for $g_i$s.
    Then the $f_i$ and $g_i$ glue uniquely together to give $f\colon W\to X$ and $g\colon W\to Y$, and
    because locally we have $\alpha f=\beta g$ we must have this globally.
    \qed

\eproof

\bdefn

    A subfunctor $h'\inj h$ is an {\emphcolor open subfunctor} if for all schemes $X$ and maps
    $h_X\to h$, the fibered product $h_X\times_hh'$ is representable, say by a scheme $U$, and
    $h_U\to h_X$ (induced from $h_X\times_hh'\to h_X$) corresponds to an open embedding of schemes
    $U\to X$.

\edefn

\blemm

    A morphism $h_W\to h_Z$ is an open subfunctor iff $W\to Z$ is an open embedding.

\elemm

\bproof

    Suppose that $h_W\to h_Z$ is an open subfunctor, then taking $X=Z$ we have that
    $h_Z\times_{h_Z}h_W\cong h_W$ (which is already representable) and then by definition $W\to Z$ is an
    open embedding.

    Conversely if $W\to Z$ is an open embedding then for a scheme $X$ by \refmath[lemma]{pullbackopen}
    we know that $X\times_ZW$ exists and is isomorphic to $h_U$ where $U\inj X$ is open.
    Thus $h_X\times_{h_Z}h_W$ is representable by $h_U$ and $U\inj X$ is an open embedding.
    \qed

\eproof

\blemm

    If $h''\inj h'$ and $h'\inj h$ are open subfunctors, then $h''\inj h$ is an open subfunctor.

\elemm

\bproof

    Let $X$ be a scheme, and so
    $$ h_X\times_hh'' \cong h_X\times_h(h'\times_{h'}h'') \cong (h_X\times_hh')\times_{h'}h'' . $$
    By assumption there is an open subscheme $U\inj X$ so that $h_X\times_hh'\cong h_U$, so
    $$ h_X\times_hh'' \cong h_U\times_{h'}h'' , $$
    and again there is an open subscheme $V\inj U$ such that this is isomorphic to $h_V$, and
    $V\inj U\inj X$ is an open embedding as desired.
    \qed

\eproof

\blemm

    \benum
        \item If $h'\inj h$ and $h''\inj h$ are open subfunctors, their pullback $h'\times_hh''$
        is an open subfunctor of $h',h'',h$.
        We call the pullback their {\emphcolor intersection}.
        \item If $U,V\inj X$ are open subschemes of $X$, then the intersection of $h_U,h_V\inj h_X$
        is $h_{U\cap V}$.
    \eenum

\elemm

\bproof

    For (1), let $X$ be a scheme so that
    $$ h_X\times_{h'}(h'\times_hh'') \cong h_X\times_hh'' \cong h_U $$
    for some open subscheme $U\inj X$.
    Thus $h'\times_hh''$ is an open subfunctor of $h'$, and by symmetry for $h''$.
    As these are open subfunctors of $h$, by the previous lemma the intersection is an open subfunctor of $h$.

    For (2) we need to show that $h_U\times_{h_X}h_V\cong h_{U\cap V}$.
    This follows directly from \refmath[lemma]{pullbackopen}.
    \qed

\eproof

\blemm[name={fiberopenisopen}]

    Let $\alpha\colon X\to Z$ and $\beta\colon Y\to Z$ be morphisms of schemes.
    Let $U\subseteq X,V\subseteq Y,W\subseteq Z$ be open subschemes so that $\alpha(U),\beta(V)\subseteq W$.
    Show that the natural map $h_U\times_{h_W}h_V\inj h_X\times_{h_Z}h_Y$ is an open embedding.

\elemm

\bproof

    We want to show that for any $h_T\colon h_X\times_{h_Z}h_Y$, the pullback
    $h_T\times_{h_X\times_{h_Z}h_Y}h_V\times_{h_W}h_U$ is represented by an open subscheme of $T$.
    The map $h_T\to h_X\times_{h_Z}h_Y$ is determined by morphisms $\phi\colon T\to X$ and
    $\psi\colon T\to Y$.
    Then we claim that
    $$ h_T\times_{h_X\times_{h_Z}h_Y}h_V\times_{h_W}h_U \cong h_{\phi^{-1}U\cap\psi^{-1}V} . $$
    Indeed a pair of morphisms $h\to h_T$ and $h\colon h_V\times_{h_W}h_U$ gives morphisms
    from $h$ to $h_T,h_V,h_U$ which commute with the base morphisms, and thus a morphism
    $h\to h_U\times_{h_X}h_T$ and $h\to h_V\times_{h_Y}h_T$.
    Because $U\subseteq X$ and $V\subseteq Y$ are open, we know that these fiber products are
    $h_{\phi^{-1}U}$ and $h_{\psi^{-1}V}$, so we have morphisms $h\to h_{\phi^{-1}U}$ and
    $h\to h_{\psi^{-1}V}$.
    These commute with the base morphisms giving us a morphism
    $h\to h_{\phi^{-1}U}\times_{h_T}h_{\psi^{-1}V}=h_{\phi^{-1}U\cap\psi^{-1}V}$.
    All of the steps here gave unique morphisms, so we have shown the desired result.
    \qed

\eproof

\bdefn

    Let $h_i$ be a collection of open subfunctors of $h$.
    Then for a scheme $X$ and a morphism $h_X\to h$, we have that $h_X\times_hh_i=h_{U_i}$ for some open
    subscheme $U_i\inj X$.
    Then $h_i$ {\emphcolor covers} $h$ if the $U_i$ cover $X$ for all $X$.

\edefn

\blemm[name={pullbackcover}]

    Suppose we have morphisms of schemes $X\to Z$ and $Y\to Z$.
    Let $\set{Z_i}_i$ be an affine cover of $Z$, and $\set{X_{ij}}_j$ be an affine cover of $Z_i$'s
    preimage in $X$, and $\set{Y_{ij}}_j$ an affine cover of $Z_i$'s preimage in $Y$.
    Then $\set{h_{X_{ij}}\times_{h_{Z_i}}h_{Y_{ik}}}_{ijk}$ is an open cover of $h_X\times_{h_Z}h_Y$.

\elemm

\bproof

    By \refmath[lemma]{fiberopenisopen} these are open subfunctors of $h_X\times_{h_Z}h_Y$.
    By the same lemma for $h_T\to h_X\times_{h_Z}h_Y$ corresponding to $T\xto\phi X$ and $T\xto\psi Y$,
    we have that
    $$ h_T\times_{h_X\times_{h_Z}h_Y}h_{X_{ij}}\times_{h_{Z_i}}h_{Y_{ik}}\cong
    h_{\phi^{-1}X_{ij}\cap\psi^{-1}Y_{ik}} . $$
    So we claim that $\phi^{-1}X_{ij}\cap\psi^{-1}Y_{ik}$ cover $T$.
    Indeed for $t\in T$ we have that $\alpha\phi(t)=\beta\psi(t)\in Z_i$ for some $i$.
    Then $\phi(t)\in X_{ij}$ for some $j$ and $\psi(t)\in Y_{ik}$ for some $k$, and thus
    $t\in\phi^{-1}X_{ij}\cap\psi^{-1}X_{ik}$ as desired.
    \qed

\eproof

\blemm

    Let $h$ be a Zariski sheaf which has an open cover by representable functors.
    Then $h$ is representable.

\elemm

\bproof

    Let $h_{U_i}\inj h$ be an open cover.
    Then $h_{U_i}\times_hh_{U_j}$ is representable by some open $U_{ij}\inj U_i$.
    It follows then that $U_{ij}\cong U_{ji}$ and so we can glue the $U_i$ along the $U_{ij}$s to get a
    scheme $X$.

    Now we construct a morphism $h_X\to h$.
    By Yoneda this is equivalent to giving an element of $h(X)$.
    Because $h$ is a Zariski sheaf, we can do this locally.
    The $U_i$ cover $X$, so we can give compatible sections of $h(U_i)$, which is the same as
    $h_{U_i}\to h$.
    Taking the given open embeddings $h_{U_i}\inj h$ we see that we may glue them together (they agree
    on overlaps) to give us $h_X\to h$.

    Now let $Y$ be another scheme, we want to show that $h_X(Y)\to h(Y)$ is an isomorphism.
    Given a morphism $h_Y\to h$ we have that $h_Y\times_hh_{U_i}\cong h_{V_i}$ cover $Y$.
    Then we have morphisms $h_{V_i}\to h_{U_i}\inj h_X$, which glue together to give $h_Y\to h_X$.
    And conversely if $h_Y\to h_X$ is a morphism, then this gives us a morphism $h_Y\to h_X\to h$.
    These two procedures are natural inverses, so we see that $\hom(\bul,h_X)\cong\hom(\bul,h)$
    and thus $h\cong h_X$ as desired.
    \qed

\eproof

\bthrm

    The fiber products of schemes exists.

\ethrm

\bproof

    We know that $h_X\times_{h_Z}h_Y$ is a Zariski sheaf by \refmath[lemma]{fiberzariski}.
    Furthermore by \refmath[lemma]{pullbackcover} we can cover this by representable functors.
    Thus by the previous lemma $h_X\times_{h_Z}h_Y$ is representable, and so $X\times_ZY$ exists.
    \qed

\eproof

\subsubsection{Computing fiber products}

There are a few ways of computing fiber products (base change).
The first is if we want to base change relative to an open embedding, then we know what the fiber product
looks like by \refmath[lemma]{pullbackopen}.

Secondly we know that for rings $A\otimes_BB[x_1,\dots,x_n]\cong A[x_1,\dots,x_n]$.
We form morphisms in both directions by mapping $a\otimes p(\bar x)\mapsto ap(\bar x)$, and conversely
mapping $x_i$ to $1\otimes x_i$ and using the universal property of polynomial algebras.
In particular $A[\bar x]\cong A\otimes_\bZ\bZ[\bar x]$.

Thus we have that
$$ \spec A\times_{\spec B}\bA^n_B \cong \bA^n_A . $$
In particular
$$ \bA^n_A \cong \spec A\times_{\spec\bZ}\bA^n_\bZ . $$
In general then we define the following.

\bdefn

    Let $X$ be a scheme, then we define $\bA^n_X$ by $X\times_{\bZ}\bA^n_\bZ$.

\edefn



